\documentclass[english,twoside, openright, hidelinks, 11 pt, class=report,crop=false]{standalone}
\input{../../preamb}
\input{../bm}


\begin{document}
\section{Midtpunktsmetoden}
Gitt en kontinuerlig funksjon $f(x)$, hvor $f(a_0)$ og $f(b_0)$ har forskjellige fortegn. For å finne et tilnærmet nullpunkt $x_t$ til $f$, kan vi anvende \outl{halveringsmetoden}\index{halveringsmetoden}. Tre iterasjoner ved denne metoden kan beskrives slik:
\st{
\textsl{Figur a}
\begin{itemize}
	\item Vi setter $m_0=\frac{a_0+b_0}{2}$.
	\item Vi sjekker om ${f(m_0)=0}$ (i så fall er ${x_t=m_0}$). Dette er ikke tilfellet, og dermed går vi videre.
\end{itemize}
\textsl{Figur b}
\begin{itemize}
	\item Vi sjekker om $f(a_0)$ og $f(m_0)$ har ulikt fortegn. I så fall er $f(a_0)\cdot f(m_0)<0$. Hvis ulikheten er sann, er $x_t\in[a_0, m_0]$.  Hvis ulikheten er usann, er $x_t\in[m_0, b_0]$. Ulikheten er usann, og dermed setter vi $a_1=m_0$ og $b_1=b_0$.
	\item Vi setter $m_1=\frac{a_1+b_1}{2}$.
	\item Vi sjekker om ${f(m_1)=0}$. Dette er ikke tilfellet, og dermed går vi videre.
\end{itemize}
\textsl{Figur c}
\begin{itemize}
	\item Vi sjekker om ${f(a_0)\cdot f(m_0)<0}$. Ulikheten er sann, og dermed setter vi $a_2=a_1$ og $b_2=m_1$.
	\item Vi setter $m_2=\frac{a_2+b_2}{2}$.
	\item Vi sjekker om ${f(m_2)=0}$. Dette er ikke tilfellet, og dermed går vi videre (med mindre vi vurderer $m_2$ som en god nok tilnærming).
	
\end{itemize}
\begin{figure}
	\centering
	\subfloat[]{\includegraphics{\figp{bisectmeth_a}}}\;
	\subfloat[]{\includegraphics{\figp{bisectmeth_b}}}\;
	\subfloat[]{\includegraphics{\figp{bisectmeth_c}}} 
	%\subfloat[]{\includegraphics{\figp{bisectmeth_d}}}
\end{figure}
}


\reg[Halveringsmetoden]{
Gitt en kontiunerlig funksjon $f(x)$, hvor $f(a_0)$ og $f(b_0)$ har forskjellige fortegn. La $x_n$ for ${n\in\mathbb{N}}$ være en tilnærming til et nullpunkt til $f$ på intervallet $[a_0, b_0]$. Ved halveringsmetoden er $x_n$ gitt ved
\[ x_n=\frac{a_n+b_n}{2} \]
For $n>0$ har vi at
\begin{itemize}
	\item Hvis $f(a_{n-1})\cdot f(m_{n-1})<0$, er $a_n=a_{n-1}$ og $b_n=m_{n-1}$. Hvis ikke, er $a_n=m_{n-1}$ og $b_n=b_{n-1}$.
\end{itemize}
}
\spr{
\textit{Halveringsmetoden} kalles også \outl{midtpunktsmetoden}.
}

\section{Newtons metode}
Gitt en funskjon $ f(x) $ si at vi ønsker å finne et tall $ a $ slik at $ {f(a)=0} $. Ved \outl{Newtons metode} gjør vi denne antakelsen for å en tilnærming $ a $: \regv
\st{
La $ x_1 $ være skjæringspunktet mellom $ x $-aksen og tangenten til $ f $ i $ x_0 $. Vi antar da at $ |x_1-a|<|x_0-a| $. Sagt med ord antar vi at $ x_1 $ gir en bedre tilnærming for $ a $ enn det $ x_0 $ gjør.
}
Siden $ x_1 $ er skjæringspunktet mellom $ x $-aksen og tangenten til $ f $ i $ x_0 $, har vi  at\footnote{Se oppgave??}
\alg{
f'(x_0)(x_1-x_0)+f(x_0)&=0 \\
f'(x_0)x_1 &= f'(x_0)x_0-f(x_0) \\
x_1 &= x_0-\frac{f(x_0)}{f'(x_0)}
}
\begin{figure}
\subfloat[]{
\includegraphics{\figp{newt1a}}
}
\subfloat[]{
	\includegraphics{\figp{newt1b}}
}
\end{figure}
La $ x_2 $ være skjæringspunktet mellom $ x $-aksen og tangenten til $ f $ i $ x_1 $. Ved å gjenta prosedyren vi brukte for å finne $ x_1 $, kan vi finne $ x_2 $, som vi antar er en enda bedre tilnærming for $ a $ enn $ x_1 $.
Prosedyren kan vi gjenta fram til vi har funnet en $ x $-verdi som gir en tilstrekkelig\footnote{Hva som er en \textit{tilstrekkelig tilnærming} er det opp til oss selv å bestemme.} tilnærming til $ a$.
\newpage
\reg[Newtons metode]{
Gitt en funskjon $ f(x) $. Si at vi ønsker å finne et tall $ a $ slik at $ {f(a)=0} $. Gitt $ x $-verdiene $ x_{n} $ og $ x_{n+1} $ for $ {n\in\mathbb{N}} $. Ved å bruke formelen
\nn{
x_{n+1} = x_n-\frac{f(x_n)}{f'(x_n)} 
}
antas det at $ x_{n+1} $ gir en bedre tilnærming for $ a $ enn $ x_{n} $.
}
\spr{
\outl{Newtons metode} kalles også \outl{Newton-Rhapsos metode}.
}
\info{Når er tilmærmingen god nok?}{
Newtons metode beskriver en iterasjonsprossess som man håper at nærmer seg en verdi. Hvis meotden lykkes, vil $ x_{n+1}$ og $ x_n $ etterhvert være veldig like, og slik kan en grense for hvor liten $ {|x_{n+1}-x_n|} $ kan være fungere som et godt mål for når iterasjonsprossessen skal stoppe.
}
\section{Trapesmetoden} \label{Trapesmetoden}
Gitt en funksjone $ f(x) $. Integralet $ \int_a^b f \,dx $ kan vi tilnærme ved å 
\st{
\begin{enumerate}
	\item Dele intervallet $ [a, b] $ inn i mindre intervall. Disse kaller vi \outl{delintervall}.
	\item Finne en tilnærmet verdi for integralet av $ f $ på hvert\\ delintervall.
	\item Summere verdiene fra punkt 2.
\end{enumerate}
}\regv

I \fref[a]{trapmetfig} har vi 3 like store delintervaller. Hvis vi setter $ {a=x_0} $ og $ {\Delta x=\frac{b-a}{3}} $, betyr dette at 
\alg{
x_1&=x_0+\Delta x & x_2&=x_0+2\Delta x & x_3&=x_0+3\Delta x=b
}
En tilnæret verdi for $ \int_{a}^{x_1}f\,dx $ får vi ved å finne arealet til trapeset med hjørner (husk at $ x_0=a $)
\alg{
(x_0, 0) && (x_1, 0) && (x_1, f(x_1)) && (x_0, f(x_0))
}
Dette arealet er gitt ved uttrykket
\[ \frac{1}{2}(x_1-x_0)\left[f(x_0)+f(x_1)\right]=
\frac{\Delta x}{2}\left[f(x_0)+f(x_1)\right] \]
Ved å tilnærme integralet for hvert delintervall på denne måten, kan vi skrive
\[ \int_{a}^{b} f\,dx \approx \frac{\Delta x}{2}\sum_{i=0}^{2} \left[f(x_i)+f(x_{i+1})\right] \]
\begin{figure}
	\centering
	\subfloat[Tilnærming med 3 del-intervaller.]{\includegraphics{\figp{trapmet}}}\qquad
\subfloat[Tilnærming med 20 del-intervaller]{\includegraphics{\figp{trapmetb}}}
	\caption{\label{trapmetfig}}
\end{figure}

\reg[Trapesmetoden \label{trapmet}]{
Gitt en integrerbar funksjon $ f $. En tilnærmet verdi for $ \int_{a}^{b} f\,dx $ er da gitt som
\begin{equation}\label{trapmeteq}
	\int_{a}^{b} f\,dx \approx \frac{\Delta x}{2}\sum_{i=0}^{n-1}\left[f(x_i)+f(x_{i+1})\right]
\end{equation}
hvor
\algv{
n&\in \mathbb{\hat{N}}\vn
a&=x_0 \vn
b&=x_n \vn
\Delta x &= \frac{b-a}{n}\vn
x_{n+1}&=x_0+i\Delta x
}
}
\info{\note}{
Slik \rref{trapmet} er formulert, vil $ {[a, b]} $ være delt inn i $ {n} $ delintervaller.
}

\end{document}